%%%%%
%
\section{Conversion du format \texorpdfstring{\texttt{double}}{double}
vers {\texorpdfstring{\texttt{float}}{float}} }

\itemstyle{Op\'erateur~:} \operfmt{CnvDbl2Flt}
\index{CnvDbl2Flt}

%%%%%OBJET%%%%
\itemstyle{Objet~:}
Transcodage d'un param\`etre du format \texttt{double} vers le format \texttt{float}.

%%%%%ARGUMENTS%%%%
\bigskip
%%%%%ENTREES%%%%
\begin{tabarg}
\multicolumn{3}{l}{\titretab {Entr\'ees}}\\
\hline
$X_d$ & Vecteur, \texttt{double} & Param\`etre \`a convertir (unit\'e indiff\'erente)\\
\end{tabarg}

\medskip
%%%%%SORTIES%%%%
\begin{tabarg}
\multicolumn{3}{l}{\titretab {Sorties}}\\
\hline
$X_f$ & Vecteur, \texttt{float} & $X_d$ converti en \texttt{float} (m\^eme unit\'e que $X_d$)\\
\end{tabarg}


%%%%%FORMULE%%%%
\itemstyle{Formule}

%%CnvDbl2Flt
\[
X_f = X_d
\]



%%%%%COMMENTAIRE%%%%
\itemstyle{Commentaire~:}
Cette op�ration comporte des risques de perte de pr�cision, et 
de d�bordement. C'est � l'utilisateur de s'assurer que la
pr�cision r�elle du param�tre d'entr�e correspond � la 
pr�cision d'un \texttt{float}, et qu'il n'y a pas de
valeurs sup�rieures (en valeur absolue) � \texttt{FLT\_MAX}.

%%%%%AUTEUR%%%%
\itemstyle{Auteur~:}
CNRM/GMEI/TRAMM

%%%%%REFERENCE%%%%
\itemstyle{R\'ef\'erences~:}
